\setcounter{section}{60} % tema número N-1
\section{Sistemas Microsoft}

Nombre completo: Sistemas operativos Microsoft. Fundamentos, Administración, Instalación, Gestión.

\localtableofcontents

\subsection{Conceptos}
    La arquitectura se basa en dos niveles, el nivel usuario y el núcleo. 

    A nivel usuario los recursos están limitados en su acceso. En el entorno se trabaja con una interfaz entre las aplicaciones y el sistema. Estas interfaces son bibliotecas dinámicas. En un principio NT admitía tres interfaces> OS/2, POSIX y Win32. La última tiene llamadas para crear procesos e hilos y asegurar seguridad.

    Windows proporciona entornos de ejecución especiales llamados \texttt{Windows Over Windows}. WoW32 se usa en los sistemas x86 para ejecutar aplicaciones de Windows 3.x, WoW 64 se usa para ejecutar aplicaciones de 32bits en sistemas x64.

    A nivel de núcleo se tiene acceso a memoria y dispositivos externos. Se ejecuta en un área protegida de la memoria. Ahí se controla el acceso a planificación, gestión de hilos, de memoria.

    \subsubsection{Hardware Abstraction Layer}
        Esta capa presenta una interfaz, ocultando los detalles técnicos específicos.

    \subsubsection{Drivers}
        Son un componente software que permite al sistema comunicarse con los dispositivos. Son bibliotecas dinámicas. El administrador de E/S organiza una ruta de flujo de datos para cada instancia. Los drivers se ejecutan en modo kernel.

    \subsubsection{MicroKernel}
        Por encima del nivel de abstracción hardware está el \texttt{NTOS} con kernel y ejecutivo.

        Kernel puede referirse a todo el código que se ejecuta en modo kernel, pero también al ntoskrnl.exe que contiene al código, y también a la capa de kernel dentro.

        Proporciona abstracciones para la CPU. incluyendo la de los hilos, excepciones, traps e interrupciones.

        Para los hilos, el kernel es responsable de determina qué se ejecuta, hasta que se termine su quantum o tiene que esperar a que ocurra algo.

    \subsubsection{Ejecutivo}
        Se relaciona con todos los subsistemas en modo usuario, se encarga de la entrada/salida, seguridad, gestión de objetos, de procesos y de caché.
\subsection{Conceptos básicos}

    \subsubsection{Procesos}
        Windows trabaja con procesos, que son instancias en ejecución de una aplicación. Pueden existir varias copias de un mismo proceso con su propio espacio de memoria y recursos.

        Algunos especiales son \texttt{SYSTEM} para el kernel, \texttt{SMSS.EXE} para administrar sesiones, \texttt{WINLOGON.EXE} para las conexiones y sesiones de usuario, \texttt{CSRSS.EXE} para el entorno win32.

        Los procesos usan hilos para ejecutar código (cada proceso tiene que tener al menos un hilo). Un caso son las fibras, hilos especiales con su propia pila y stack de instrucciones, no se gestionan por el SO sino por el usuario.

        Una Tarea es un conjunto de hilos creados por un proceso. La fibra y la Tarea se usan pocas veces, normalmente quedándose en Procesos e Hilos.

    \subsubsection{Registro de windows}
        El \texttt{Registro de windows} se usa para gestionar a bajo nivel el sistema. Se organiza en volúmenes separados. Cada uno de estos volúmenes tiene una estructúra jerárquica con directorios y pares clave/valor.

    \subsubsection{Sistema de ficheros}
        Dentro del sistema de ficheros. Los dispositivos pueden estar formateados en FAT32, NTFS ReFS y exFAT.
        La estructura está descrita en el Master File Table donde se describen los registros de archivos y directorios. Microsoft desarrolló bitLocker para cifrar discos completos.

    \subsubsection{Seguridad}
        Windows no se diseñó específicamente para cumplir con requisitos de seguridad. Hereda propiedades de seguridad del diseño original de NT como:
        \begin{description}
            \item[Inicio seguro]
            \item[Controles de acceso discrecionales]
            \item[Controles de acceso privilegiados]
            \item[Protección del espacio de direcciones por proceso]
            \item[Limpieza de memoria antes de reasignarla]
            \item[Auditoría de seguridad]
            \item[Identificadores universales de usuario]
            \item[Discrecional ACL]
            \item[Privilegios]
            \item[Descriptor de seguridad]
            \item[Control de usuarios] 
        \end{description}

    \subsubsection{Intérprete}
        Windows utiliza Powershell, desarrollado en .NET, permitiendo la automatización de tareas y de configuración de productos de Microsoft.


\subsection{Versiones de Windows}
    Microsoft gestiona en paralelo versiones de escritorio y servidor. Comparten parte del desarrollo. Aunque no tienen por qué ser equivalencias completas.

    \begin{table}[H]
    \begin{tabular}{ll}
        Escritorio & Servidor \\
        \midrule
        XP         & Server 2003 \\
        Vista      & Server 2008 \\
        7          & Server 2008 R2 \\
        8          & Server 2012 \\
        8.1        & Server 2012 R2 \\
        10         & Server 2016 y Server 2022 \\
        11         & Server 2025 \\
    \end{tabular}
    \end{table}
    \subsubsection{Windows 10}
        Se introdujo una arquitectura de aplicaciones universales para poder ejecutarse en todas las familias de productos, desde móviles, ordenadores, consolas, embebidos y otros.

        Se introdujo Edge como navegador web, el acceso por huella o webcam.

    \subsubsection{Windows 11}
        Requiere un TPM para ejecutarse. Incorpora Windows Hello para iniciar sesión y ejecución de aplicaciones android con WSA.

    \subsubsection{Windows Server 2012}
        Se parte de la instalación Server Core donde se pueden ir añadiendo elementos.

        Se añade Hyper-V como elemento de virtualización, incluido redes, almacenaiento, copias de seguridad, etc.

        También se añade ReFS como sistema de archivos resiliente.

    \subsubsection{Windows Server 2012 R2}
        La interfaz gráfica pasa a ser opcional. Se permite la creación de grupos de servidores para gestionar múltiples máquinas en conjunto. También con el almacenamiento.

        Se permite replicar máquinas con Hyper-V siempre que haya una conexión de red.

        DirectAccess permite crear túneles VPN sin necesitar un servidor dedicado.

    \subsubsection{Windows Server 2016}
        Se incluye el uso de contenedores por Hyper-V. Se pueden blindar las máquinas virtuales y se permite generar guardianes de dispositivos, y otras herramientas de seguridad.

        También se permiten migrar discos duros a discos virtuales.

    \subsubsection{Windows Server 2019}
        Se integra mas con Azure. Windows Defender ATP, subsistemas de Windows para Linux.

        Se incluyen mejoras en los contenedores, disminuyendo el tamaño, incluido integración con K8 y su gestión.

    \subsubsection{Windows Server 2022}
        Secured-core server protege el núcleo frente a ataques. Servidores con Azure Arc Habilitados permiten gestión de máquinas virtuales.
    \subsubsection{Windows Server 2025}
        Incorpora sudo. y parece que migra al Sistema como servicio.